% NOTE TO ANDREW: decide whether/how to acknowledge AI assistance (Claude) before sending;
% some venues require disclosure.
%
% Build:  tectonic paper/linear-subsequence-note.tex
\documentclass[10pt,reqno]{amsart}

\usepackage[T1]{fontenc}
\usepackage[utf8]{inputenc}
\usepackage{amsmath,amssymb,amsthm}
\usepackage{booktabs}
\usepackage{array}
\usepackage{microtype}
\usepackage[margin=0.8in]{geometry}
\usepackage[colorlinks=true,linkcolor=black,citecolor=black,urlcolor=blue]{hyperref}

\theoremstyle{plain}
\newtheorem{theorem}{Theorem}
\newtheorem{lemma}[theorem]{Lemma}
\newtheorem{corollary}[theorem]{Corollary}
\newtheorem{proposition}[theorem]{Proposition}

\theoremstyle{definition}
\newtheorem{definition}[theorem]{Definition}
\newtheorem{observation}[theorem]{Observation}
\newtheorem{problem}{Open Problem}

\theoremstyle{remark}
\newtheorem{remark}[theorem]{Remark}

\numberwithin{theorem}{section}

% compact vertical spacing
\setlength{\abovedisplayskip}{5pt}
\setlength{\belowdisplayskip}{5pt}
\setlength{\abovedisplayshortskip}{3pt}
\setlength{\belowdisplayshortskip}{3pt}
\makeatletter
\def\thm@space@setup{\thm@preskip=5pt \thm@postskip=5pt}
\makeatother

\DeclareMathOperator{\sco}{sc}
\newcommand{\Zm}{\mathbb{Z}/m\mathbb{Z}}
\newcommand{\N}{\mathbb{N}}
\newcommand{\Sig}{\Sigma_k}
\newcommand{\tm}{t_{2,m}}

\begin{document}

\title[A tight family for linear subsequences]{A tight family for the state complexity of
linear subsequences of automatic sequences}

\author{Andrew Hingston}
\email{\texttt{hingstonandrew@icloud.com}}

\date{\today}

\subjclass[2020]{68Q45, 68R15, 11B85}
\keywords{automatic sequence, state complexity, factor complexity, generalised
Thue--Morse word, linear subsequence}

\begin{abstract}
Open Problem~1 of Moradi, Rampersad and Shallit (arXiv:2512.10017v5, \S4) asks for a class
of automatic sequences $(h_m(i))_{i \ge 0}$, each generated by an $m$-state msd-first DFAO,
whose linear subsequences $(h_m(ni))_{i \ge 0}$ have state complexity at least $Cnm^2$ for
infinitely many $n$ and some constant $C$. We answer it with the generalised Thue--Morse
family $h_m(i) = s_2(i) \bmod m$ ($s_2$ the binary digit sum): for every $m \ge 2$ and odd
$n \ge 3$ the minimal msd-first DFAO of $(h_m(ni))_{i \ge 0}$ has exactly
$m^2(n-1) - m\max(2^{\,p-1}, r)$ states, where $n - 1 = 2^p + r$, $0 < r \le 2^p$.
This is at least $\tfrac12 n m^2$ for \emph{every} odd $n \ge 3$, and
$\limsup_{n} \sco(h_m(ni))/(nm^2)$ equals $1 - 1/(3m) \ge 5/6$; the value $C = 1/2$ is attained
exactly at $(m,n) = (2,3)$, so within this family it is not improvable uniformly over odd
$n$. The proof joins three ingredients: a Myhill--Nerode lemma for every digit-additive
sequence and every $n$, generalising the authors' own optimality theorem; a reachability
lemma for odd $n$ from Gelfond's 1968 equidistribution theorem, replacing their
Thue--Morse-specific one; and the known factor complexity of the generalised Thue--Morse
word.
\end{abstract}

\maketitle

\section{Background and the problem}
\label{sec:bg}

Fix a base $k \ge 2$ and write $\Sig = \{0,\dots,k-1\}$; for $x \in \Sig^*$ let $[x]_k$ be
the integer it represents, most-significant digit first, and for $i \in \N$ let $(i)_k$ be
its leading-zero-free representation. A sequence $h = (h(i))_{i \ge 0}$ over a finite
alphabet $\Delta$ is \emph{$k$-automatic} if $h(i) = \tau(\delta(q_0,(i)_k))$ for some DFAO
$M = (Q,\Sig,\delta,q_0,\Delta,\tau)$. Input is msd-first with leading zeros allowed
throughout, so we assume $M$ is \emph{zero-consistent}: $\delta(q_0,0) = q_0$, which is
exactly what makes the recursion $h'(ki+b) = \delta(h'(i),b)$ valid at $i = 0$ as well. The
\emph{interior sequence} is $h'(i) = \delta(q_0,(i)_k) \in Q$, generated by the
\emph{intrinsic} DFAO $\widehat{M} = (Q,\Sig,\delta,q_0,Q,\mathrm{id})$. We write
$\rho_x(n)$ for the number of length-$n$ factors of $x$ and $\sco(x)$ for the number of
states of the minimal msd-first DFAO generating $x$.

Moradi, Rampersad and Shallit \cite{MRS} relate these two quantities. For $n \ge 1$ and
$0 \le c < n$, their Theorem~9(a) generates $(h(ni+c))_{i \ge 0}$ by a DFAO whose states are
length-$n$ factors of $h'$, the state reached on input $x$ being the window
$\bigl[h'(n[x]_k),\, \dots,\, h'(n[x]_k+n-1)\bigr]$ with output
$[r_0,\dots,r_{n-1}] \mapsto \tau(r_c)$. Hence $\sco(h(ni+c)) \le \rho_{h'}(n)$; with
$\rho_x(n) \le kn\,\rho_{\widehat M}(2) \le knm^2$ \cite[Thm.~10.3.1]{AS} this gives their
Corollary~10(a): $\sco(h(ni+c)) \le knm^2$. They then write \cite[\S4]{MRS}:

\begin{quote}
We do not know if the bounds in Corollary 10 are tight. This seems worthy of further
study, so we state it as an open problem.
\end{quote}

\begin{problem}[{\cite[\S4]{MRS}}]
\label{op1}
Find a class of automatic sequences $(h_m(i))_{i\ge0}$, generated by a DFAO of $m$ states
with msd-first input, such that there is a constant $C$ so the linear subsequences
$(h_m(ni))_{i\ge0}$ have at least the state complexity $Cnm^2$ for infinitely many $n$.
\end{problem}

(Quoted verbatim from arXiv:2512.10017v5, 2 April 2026. It appears again, verbatim, as
Problem~1 of Moradi's MMath thesis \cite{Moradi-thesis}, whose ``More Open Problems''
section leaves it unresolved, so the problem is open as of its author's latest treatment.)

For $m = 2$ the authors settle the state count: their Theorem~16 shows the Theorem~9
construction optimal for the Thue--Morse word $t$ and odd $n$; their Lemma~15, quoting
\cite{BSCRF}, shows every length-$n$ factor of $t$ occurs at a position divisible by odd
$n$; and with Brlek's factor complexity of $t$ \cite{Brlek} this yields
$\sco(t(ni)) = \rho_t(n/\nu_2(n))$ (their Theorem~14) and the explicit Corollary~17. What is
missing is the $m$-dependence --- a family in which the $m^2$ of Corollary~10 is visible.

\section{The family}
\label{sec:family}

For $m \ge 2$ let $s_2(i)$ be the number of $1$s in the binary representation of $i$, and
let
\[
  \tm = \bigl(s_2(i) \bmod m\bigr)_{i \ge 0}
\]
be the \emph{generalised Thue--Morse word} over $\Zm$; for $m = 2$ this is the Thue--Morse
word. It is generated by the $m$-state zero-consistent msd-first DFAO
\[
  M_m = \bigl(\Zm,\ \{0,1\},\ \delta,\ 0,\ \Zm,\ \mathrm{id}\bigr),
  \qquad \delta(q,d) = q + d \bmod m ,
\]
equivalently by the $2$-uniform morphism $j \mapsto j\,(j{+}1)$ on $\Zm$ with identity
coding. State $q$ is reached by $1^q$ and the coding is injective, so distinct states are
separated by the empty word: $M_m$ is minimal, $\sco(\tm) = m$, and the interior sequence of
$M_m$ is $\tm$ itself. So $\tm$ is a legitimate candidate for Problem~\ref{op1}. The one
structural property we use is that it is \emph{digit-additive}.

\begin{definition}
A sequence $h : \N \to \Zm$ is \emph{$k$-digit-additive} if
\begin{equation}
  h(k^L X + y) \equiv h(X) + h(y) \pmod m
  \qquad\text{for all } X \ge 0,\ L \ge 0,\ 0 \le y < k^L .
  \label{eq:additive}
\end{equation}
\end{definition}

Equivalently, $h(i) = \sum_j w_{d_j} \bmod m$ is a weighted digit sum with weights
$w : \Sig \to \Zm$ satisfying $w_0 = 0$: the representation of $k^L X + y$ is that of $X$
followed by the zero-padded length-$L$ representation of $y$, and $w_0 = 0$ makes the
padding harmless. Condition \eqref{eq:additive} forces $h(0) = 0$, and the associated DFAO
$\delta(q,d) = q + w_d$ is zero-consistent. Here $\tm$ is the case $k = 2$,
$(w_0,w_1) = (0,1)$.

\section{Lemma A: Nerode classes are exactly the windows}
\label{sec:lemA}

\begin{lemma}[Lemma A]
\label{lemA}
Let $k \ge 2$, $m \ge 2$, let $h : \N \to \Zm$ be $k$-digit-additive, and let $n \ge 1$.
For $j \ge 0$ put
\[
  W(j) \;=\; \bigl(h(nj + g)\bigr)_{0 \le g < n} \;\in\; (\Zm)^n .
\]
Then two words $x,y \in \Sig^*$ are Myhill--Nerode equivalent for the sequence
$(h(ni))_{i \ge 0}$ if and only if $W([x]_k) = W([y]_k)$. Consequently
\[
  \sco\bigl(h(ni)\bigr) \;=\; \#\{\,W(j) \;:\; j \ge 0\,\},
\]
that is, the reachable part of the Theorem~9(a) automaton of \cite{MRS} is already minimal.
\end{lemma}

\begin{proof}
Recall that $x \sim y$ means $h(n[xz]_k) = h(n[yz]_k)$ for every $z \in \Sig^*$.

\emph{Sufficiency.} Let $z \in \Sig^L$, so $n[xz]_k = k^L \cdot n[x]_k + n[z]_k$. Write
$n[z]_k = k^L g + w$ with $0 \le w < k^L$ and $g = \lfloor n[z]_k / k^L \rfloor$; since
$[z]_k \le k^L - 1$ we have $n[z]_k < nk^L$, hence $g \le n-1$. Thus
$n[xz]_k = k^L(n[x]_k + g) + w$, and \eqref{eq:additive} gives
\begin{equation}
  h\bigl(n[xz]_k\bigr) \;=\; h\bigl(n[x]_k + g\bigr) + h(w)
  \;=\; W([x]_k)[g] + h(w) \pmod m .
  \label{eq:key}
\end{equation}
The right-hand side depends on $x$ only through $W([x]_k)$, so $W([x]_k) = W([y]_k)$
implies $x \sim y$.

\emph{Necessity.} Suppose $W([x]_k)[g] \ne W([y]_k)[g]$ for some $0 \le g < n$. Choose $L$
with $k^L \ge n$ and set $\varphi(Z) = \lfloor nZ/k^L \rfloor$ for $0 \le Z < k^L$. Then
$\varphi(0) = 0$, $\varphi(k^L-1) = \lfloor n - n/k^L \rfloor = n-1$ (as $0 < n/k^L \le 1$),
and $\varphi(Z+1) - \varphi(Z) \le \lceil n/k^L \rceil = 1$, so $\varphi$ is onto
$\{0,\dots,n-1\}$. Pick $Z$ with $\varphi(Z) = g$ and let $z \in \Sig^L$ be the zero-padded
representation of $Z$. Applying \eqref{eq:key} for both $x$ and $y$, the $h(w)$ coincide and
\[
  h\bigl(n[xz]_k\bigr) - h\bigl(n[yz]_k\bigr) \;=\; W([x]_k)[g] - W([y]_k)[g] \;\ne\; 0
  \pmod m ,
\]
so $x \not\sim y$.

Finally $W([x]_k)$ depends on $x$ only through $[x]_k$, and every $j \ge 0$ equals $[x]_k$
for some $x$, so the Nerode classes biject with $\{W(j) : j \ge 0\}$.
\end{proof}

\begin{remark}
Lemma~\ref{lemA} generalises \cite[Thm.~16]{MRS} (the case $k = m = 2$, $n$ odd), and the
proof is theirs. Two points are worth isolating: it holds for \emph{every} $n$, even or odd,
and it needs \emph{no} reachability hypothesis. For digit-additive sequences the
Theorem~9(a) construction is thus always optimal; the only remaining question is which
windows occur at multiples of $n$.
\end{remark}

\section{Lemma B: for odd $n$, every factor occurs at a multiple of $n$}
\label{sec:lemB}

\begin{lemma}[Lemma B]
\label{lemB}
Let $m \ge 2$, $h = \tm$, and let $n \ge 1$ be odd. Then every length-$n$ factor of $h$
occurs at some position divisible by $n$. Consequently, by Lemma~\ref{lemA},
\[
  \sco\bigl(\tm(ni)\bigr) \;=\; \rho_{\tm}(n) .
\]
\end{lemma}

\begin{proof}
Write $s = s_2$ and fix $t$ with $2^t \ge n$. Every position $P$ factors uniquely as
$P = 2^t a + b$ with $0 \le b < 2^t$. For $0 \le g < n \le 2^t$ there are two cases: if
$b + g < 2^t$ then $h(P+g) = h(a) + h(b+g)$, while if $b+g \ge 2^t$ then
$0 \le b+g-2^t < 2^t$ and $P + g = 2^t(a+1) + (b+g-2^t)$, so
$h(P+g) = h(a+1) + h(b+g-2^t)$; both by digit-additivity. Hence the length-$n$ factor at
$P$ is a function of the \emph{type}
\[
  \bigl(u,v,b\bigr) \;=\; \bigl(h(a),\ h(a+1),\ b\bigr) .
\]
It therefore suffices, given $P$ of type $(u,v,b)$, to produce $j \ge 0$ such that $nj$ has
the same type.

\emph{Step 1: available types at multiples of $n$.} As $n$ is odd it is invertible modulo
$2^t$; choose $j_0 \in [0,2^t)$ with $n j_0 \equiv b \pmod{2^t}$. Then $nj_0 - b$ is a
multiple of $2^t$ exceeding $-2^t$, so $a_0 := (nj_0 - b)/2^t \in \N$. For
$j = j_0 + 2^t j'$, $j' \ge 0$,
\[
  nj \;=\; 2^t\bigl(a_0 + n j'\bigr) + b ,
\]
so the high parts realisable with this $b$ are exactly the integers $a' \ge a_0$ with
$a' \equiv a_0 \pmod n$.

\emph{Step 2: hitting the pair $(u,v)$.} Choose $T \in \{1,\dots,m\}$ with
$T \equiv u - v + 1 \pmod m$, choose any $L > T$, and consider
\[
  a' \;=\; 2^L A + \bigl(2^T - 1\bigr), \qquad A \ge 0 .
\]
Bits $0,\dots,T-1$ of $a'$ are $1$ and bit $T$ is $0$ (here $L > T$ is used), so
$s(a') = s(A) + T$; and $a'+1 = 2^L A + 2^T$, so $s(a'+1) = s(A) + 1$, no carry reaching
$A$. Hence if $h(A) \equiv v - 1 \pmod m$ then
\[
  h(a'+1) = (v-1) + 1 = v, \qquad h(a') = (v-1) + T \equiv u \pmod m ,
\]
as required.

\emph{Step 3: the congruence.} The constraint $a' \equiv a_0 \pmod n$ reads
$2^L A \equiv a_0 - 2^T + 1 \pmod n$, and $2$ is invertible modulo the odd number $n$, so it
is a single congruence $A \equiv A_0 \pmod n$. By Gelfond's theorem \cite{Gelfond} --- for a
base $q \ge 2$ and a modulus $m$ with $\gcd(m,q-1) = 1$, the digit sum $s_q$ is
equidistributed modulo $m$ along every arithmetic progression; for $q = 2$ the hypothesis
reads $\gcd(m,1) = 1$ and is automatic --- there are infinitely many $A \equiv A_0
\pmod n$ with $s(A) \equiv v-1 \pmod m$. Choosing such an $A$ large enough that
$a' \ge a_0$ and setting $j' = (a' - a_0)/n \ge 0$, $j = j_0 + 2^t j'$, the multiple $nj$
has type $(u,v,b)$ and hence carries the same length-$n$ factor as $P$.

By Lemma~\ref{lemA}, $\sco(\tm(ni)) = \#\{W(j) : j \ge 0\}$, and the $W(j)$ are exactly the
length-$n$ factors of $\tm$ at multiples of $n$; by the above these are all of them.
\end{proof}

\begin{remark}
Lemma~\ref{lemB} is the $m$-general replacement for \cite[Lemma~15]{MRS}, which covers
$m = 2$ by citing \cite{BSCRF}; it re-proves \cite[Thm.~14]{MRS} without that citation.
Oddness of $n$ is used twice (invertibility of $n$ mod $2^t$, of $2$ mod $n$) and cannot be
dropped: for even $n$, $\#\{W(j)\} < \rho_{\tm}(n)$ strictly. This costs nothing, since
$\tm(2i) = \tm(i)$ makes $\sco(\tm(ni))$ depend only on the odd part of $n$.
\end{remark}

\section{Theorem A}
\label{sec:thmA}

The third ingredient is the factor complexity of $\tm$, determined by Tromp and Shallit
\cite{TrompShallit} and, for all bases, by Starosta \cite{Starosta}, whose Table~2 gives,
for the base-$b$ word modulo $m$ with $q$ the order of $b-1$ in $\Zm$ (so $q = m$ for
$b = 2$),
\begin{align*}
  C(n) &= qm(n-1) - m\bigl(b^{\,\kappa} - b^{\,\kappa-1}\bigr)
   && \text{for } n = b^{\,\kappa} + 1 + \ell,\ 0 \le \ell < b^{\,\kappa}-b^{\,\kappa-1},\\
  C(n) &= qm(n-1) - m\bigl(b^{\,\kappa} - b^{\,\kappa-1} + \ell\bigr)
   && \text{for } n = (2b{-}1)b^{\,\kappa-1} + 1 + \ell,\ 0 \le \ell < b^{\,\kappa+1}-2b^{\,\kappa}+b^{\,\kappa-1}.
\end{align*}
For $b = 2$ the two regimes merge.

\begin{proposition}
\label{prop:rho}
For $n \ge 3$ let $p \ge 0$ and $0 < r \le 2^p$ be the unique integers with
$n - 1 = 2^p + r$. Then $\rho_{\tm}(n) = m^2(n-1) - m\max\bigl(2^{\,p-1}, r\bigr)$,
where $2^{\,p-1}$ is read as the real number $1/2$ when $p = 0$.
\end{proposition}

The reduction is direct: the first regime is $r = \ell < 2^{\,\kappa-1}$ (with $\ell = 0$
giving $p = \kappa-1$, $r = 2^{\,\kappa-1}$), the second $r = 2^{\,\kappa-1} + \ell \ge
2^{\,\kappa-1}$; at $m = 2$ this specialises to Brlek's formula \cite{Brlek}, whence
\cite[Cor.~17]{MRS}. The case $p = 0$ occurs only for $n = 3$, where $\max(1/2,1) = 1$ and
$\rho_{\tm}(3) = 2m^2 - m$; readers preferring integers may treat $n = 3$ as a base case.

\begin{theorem}[Theorem A]
\label{thmA}
Let $m \ge 2$ and let $n \ge 3$ be odd. Write $n - 1 = 2^p + r$ with $0 < r \le 2^p$. Then
\[
  \sco\bigl(\tm(ni)\bigr) \;=\; \rho_{\tm}(n)
  \;=\; m^2(n-1) \;-\; m\max\bigl(2^{\,p-1},\,r\bigr) .
\]
Consequently, for every odd $n \ge 3$,
\[
  \sco\bigl(\tm(ni)\bigr) \;\ge\; \Bigl(1 - \tfrac{1}{2m}\Bigr) m^2 (n-1)
  \;\ge\; \tfrac34\, m^2(n-1) \;\ge\; \tfrac12\, n m^2 ,
\]
and along the odd values $n = 3\cdot 2^{\,p-1} + 1$ $(p \ge 2)$ the state count equals
$\bigl(1 - \tfrac{1}{3m}\bigr)m^2(n-1)$ exactly, so that
\[
  \limsup_{\substack{n \to \infty \\ n \text{ odd}}}
  \frac{\sco(\tm(ni))}{n m^2} \;=\; 1 - \frac{1}{3m} \;\ge\; \frac56 .
\]
In particular the family $(\tm)_{m \ge 2}$ answers Problem~\ref{op1} with $C = 1/2$ ---
for all odd $n \ge 3$, not merely infinitely many.
\end{theorem}

\begin{proof}
The first equality is Lemma~\ref{lemB} (which uses Lemma~\ref{lemA}), the second is
Proposition~\ref{prop:rho}.

For the bounds put $N = n-1 = 2^p + r$. Since $r > 0$, $2^p \le N$ and so
$2^{\,p-1} \le N/2$; since $r \le 2^p$, $N \ge 2r$ and so $r \le N/2$. Hence
$\max(2^{\,p-1},r) \le N/2$, with equality iff $r = 2^p$, and
$\sco \ge m^2N - mN/2 = (1 - \tfrac1{2m})m^2N \ge \tfrac34 m^2N$ as $m \ge 2$; finally
$N \ge \tfrac23 n$ for $n \ge 3$, giving $\sco \ge \tfrac12 nm^2$.

For the $\limsup$, minimise $f(p,r) = \max(2^{\,p-1},r)/(2^p+r)$ over $0 < r \le 2^p$: on
$r \le 2^{\,p-1}$ it is $2^{\,p-1}/(2^p+r)$, decreasing in $r$, and on $r \ge 2^{\,p-1}$ it
is $r/(2^p+r)$, increasing in $r$, so the minimum is $f = 1/3$ at $r = 2^{\,p-1}$. Thus
$\max(2^{\,p-1},r) \ge (n-1)/3$ always, with equality exactly when
$n = 3\cdot 2^{\,p-1}+1$ (odd precisely for $p \ge 2$). Hence
$\sco(\tm(ni)) \le (1 - \tfrac1{3m})m^2(n-1) < (1-\tfrac1{3m})nm^2$ for every odd
$n \ge 3$, with equality in the first inequality along $n = 3\cdot2^{\,p-1}+1$, where
$(n-1)/n \to 1$; and $1 - 1/(3m) \ge 5/6$ for $m \ge 2$.
\end{proof}

\begin{remark}
Against the bound actually proved in \cite{MRS}, namely $2nm^2$ for $k = 2$, the family
realises the fraction $(1 - \tfrac1{3m})/2 \in [\tfrac5{12}, \tfrac12)$: Corollary~10(a) is
tight in the shape $nm^2$, loose only in the factor $k$ (cf.\ Theorem~\ref{thmC}).
\end{remark}

\section{Verification}
\label{sec:verif}

Theorem~\ref{thmA} is a proof, conditional only on the two literature inputs
(Proposition~\ref{prop:rho} and Gelfond's theorem). It was nevertheless checked numerically
by implementations sharing no code. Two computed $\sco$ directly: a \emph{theory-free} msd
Myhill--Nerode search, in which the signature of $X$ is the raw table
$\bigl(u(k^L X + j)\bigr)_{L \le L_{\max},\, j < k^L}$ built from the oracle
$u(i) = s_2(ni) \bmod m$ alone, with transitions $X \mapsto kX + d$ (distinct signatures are
provably distinct states, so it cannot over-count); and a breadth-first construction of the
Theorem~9(a) window automaton followed by Moore minimisation. Separately, $\rho_{\tm}(n)$
was computed by suffix doubling on a $2^{22}$-symbol prefix, and $\#\{W(j)\}$ straight from
population counts, with no automaton involved. A further independent code path is Peanut,
which compiles the predicate \verb|T[n*i]=a| into a minimal DFA.

\begin{center}\footnotesize
\begin{tabular}{@{}>{\raggedright\arraybackslash}p{5.5cm}>{\raggedright\arraybackslash}p{5.9cm}>{\raggedright\arraybackslash}p{5.4cm}@{}}
\toprule
Check & Scope & Result \\
\midrule
Nerode search vs.\ window construction vs.\ Thm.~\ref{thmA}
  & $m = 2 \dots 7$, odd $n = 3 \dots 101$ (300 cells)
  & no mismatch \\
Thm.~\ref{thmA} vs.\ window construction, adversarial parameters
  & $m \in \{8,10,12,16,20\}$ $\times$ 7 values of $n$; 12 further cells up to $m = 32$
  & no mismatch \\
$\rho_{\tm}$ by suffix doubling vs.\ Prop.~\ref{prop:rho}
  & $m = 2 \dots 7$, every $n$ from $3$ to $130$
  & no mismatch, no non-convergence \\
$\#\{W(j) : j < J\}$ from population counts vs.\ $\rho_{\tm}(n)$ (Lem.~\ref{lemB})
  & slowest-converging cells, $J < 2\cdot10^{7}$
  & exact, e.g.\ $(m,n) = (7,101) \mapsto 4648$ \\
Peanut, \verb|dfa T[n*i]=a|
  & 14 cells, $m \le 7$, $n \le 101$
  & agrees with both constructions \\
subsequence $n = 3\cdot2^{\,p-1}+1$, $p = 2 \dots 6$
  & $m \in \{2,3,5,7\}$ (20 cells)
  & $\sco = (1-\tfrac1{3m})m^2(n-1)$ exactly \\
extremes of $\sco/(nm^2)$, $\sco/(m^2(n-1))$
  & the 300 cells above
  & minima $1/2$ (at $m=2,n=3$) and $3/4$ \\
\bottomrule
\end{tabular}
\end{center}

Two caveats about status. First, reading the state count of the per-letter DFA
\verb|dfa T[n*i]=a| as $\sco$ presumes each level set $\{i : \tm(ni) = a\}$ needs the full
$\sco$ states. This held in every cell run --- all $m$ letters gave the same count, equal to
$\sco$ --- but it is an empirical bridge, not a proof, so the engine figures corroborate the
two constructions rather than replace them. Second, and more importantly:

\begin{observation}[base 3; machine-verified only]
\label{obsB}
Let $g_m(i) = \bigl(\#\{\text{digits equal to } 2 \text{ in } (i)_3\}\bigr) \bmod m$, a
$3$-digit-additive sequence with weights $(0,0,1)$, generated by a minimal $m$-state
zero-consistent base-$3$ DFAO. Then $\sco(g_m(ni)) = \rho_{g_m}(n) = m^2(n-1)$ exactly
whenever $3 \nmid n$, and $\sco(g_m(ni)) = \sco(g_m(n'i))$ for $n = 3^{v}n'$ with
$3 \nmid n'$.
\end{observation}

Observation~\ref{obsB} was reproduced independently for $m = 2 \dots 6$ and
$n \in \{2,4,5,7,8,11,13,16,22,31,40,44\}$, and is \emph{not} proved here. Lemma~\ref{lemA}
applies verbatim, so $\sco(g_m(ni)) = \#\{W(j)\}$; the missing ingredients are an analogue
of Lemma~\ref{lemB} --- Gelfond's hypothesis $\gcd(m,q-1)=1$ fails for $q = 3$ and even $m$,
though the sequence here is a weighted digit sum rather than $s_3$ --- and a proof that
$\rho_{g_m}(n) = m^2(n-1)$. This family saturates $nm^2$, reaching $\sco/(nm^2) = 1 - 1/n$,
but only one third of the $knm^2$ of Corollary~10(a). It is recorded as a conjecture, not a
result.

\section{Remarks}
\label{sec:remarks}

\subsection*{A sharpening of Corollary 10(a)}
The window parametrisation of Lemma~\ref{lemB} is not special to $\tm$, and gives the
following, in which the factor $k$ of Corollary~10(a) becomes $1 + (k-1)/m$.

\begin{theorem}[Theorem C]
\label{thmC}
Let $h$ be generated by an $m$-state \emph{zero-consistent} msd-first DFAO with interior
sequence $h'$ and intrinsic DFAO $\widehat M$, let $n \ge 1$, and let $t$ be minimal with
$k^t \ge n$. Then
\[
  \rho_{h'}(n) \;\le\; \rho_{\widehat M}(2)\,(n-1) \;+\; m\bigl(k^t - n + 1\bigr)
  \;\le\; m^2(n-1) + m\bigl((k-1)n + 1\bigr),
\]
and hence $\sco(h(ni+c)) \le m^2(n-1) + m((k-1)n+1)$ for all $0 \le c < n$.
\end{theorem}

\begin{proof}
Write $P = k^t a + b$, $0 \le b < k^t$. Every position in $[P, P+n)$ is $k^t(a+e) + r$ with
$e \in \{0,1\}$, $0 \le r < k^t$, and $h'(k^t A + r) = \delta(h'(A), r_1 \cdots r_t)$ where
$r_1\cdots r_t$ is $r$ zero-padded to length $t$, since msd reading concatenates digit
blocks --- exactly where zero-consistency is needed. So the factor at $P$ is determined by
$(h'(a), h'(a+1), b)$, and by $(h'(a), b)$ alone when $b + n \le k^t$. There are
$k^t - n + 1$ values of $b$ of the latter kind, each contributing at most $m$ factors, and
$n-1$ of the former, each contributing at most the number $\rho_{\widehat M}(2)$ of realised
pairs $(h'(a),h'(a+1))$. The second inequality uses $\rho_{\widehat M}(2) \le m^2$ and
$k^t \le k(n-1) \le kn-1$, the latter by minimality of $t$.
\end{proof}

Zero-consistency is a genuine hypothesis, not a normalisation: without it $h'(0) = q_0$
fails, the recursion $h'(ki+b) = \delta(h'(i),b)$ breaks at $i = 0$, and the window
construction of \cite[Thm.~9]{MRS} --- hence Theorem~\ref{thmC} and every window count in
this note --- is invalid. Both families here are zero-consistent, as is every
digit-additive sequence. Theorem~\ref{thmC} improves $knm^2$ when $m \gg k$, but it is a
one-shot form of the standard blocking argument behind \cite[Thm.~10.3.1]{AS} and should be
regarded as folklore-grade.

\subsection*{On significance}
This note assembles three known pieces rather than proving a new hard theorem. The proof of
\cite[Thm.~16]{MRS} generalises essentially verbatim to Lemma~\ref{lemA}; their Lemma~15
needs a replacement, supplied by the short construction of Lemma~\ref{lemB} with Gelfond's
1968 theorem; and the growth rate $m^2n$ is already in \cite{TrompShallit, Starosta}. What
was absent from the source is precisely the join of the three. The answer appears correct
and complete, but should be read as a routine composition, and may well be known to
specialists in a form we have not located.

\subsection*{The constant is not improvable uniformly}
Within this family $C = 1/2$ is exactly attained, at $(m,n) = (2,3)$, where
$\sco = 6 = \tfrac12 \cdot 3 \cdot 2^2$. So $1/2$ is the largest constant valid for
\emph{all} odd $n \ge 3$ and all $m \ge 2$ here, and the stronger $1 - 1/(3m)$ of
Theorem~\ref{thmA} is a statement about a subsequence of $n$ only. Since Problem~\ref{op1}
asks only for infinitely many $n$, either constant answers it.

\subsection*{What is not settled}
Three questions remain open. (i) Observation~\ref{obsB}: the identity
$\rho_{g_m}(n) = m^2(n-1)$ and an analogue of Lemma~\ref{lemB} in base 3. (ii) The exact
maximum of $\sco(h(ni))$ over all $m$-state $k$-DFAOs. The bound $m^2(n-1)$ holds
exhaustively for all zero-consistent 2- and 3-state DFAOs at $k \in \{2,3\}$,
$n \in \{7,11,16,31\}$, but fails for $k = 4$: weights $(0,0,1,0)$ modulo $3$ give
$\sco(h(31i)) = 312 > 270 = m^2(n-1)$, still below Theorem~\ref{thmC}'s bound of $372$. The
truth thus lies between $m^2(n-1)$ and $m^2(n-1) + m((k-1)n+1)$, and the factor $k$ in
$knm^2$ is a real but bounded effect that vanishes as $m$ grows. (iii) The case $c > 0$,
i.e.\ $(h(ni+c))_{i\ge0}$ and \cite[Cor.~10(b)]{MRS}: only $c = 0$ is treated here, that
being what Problem~\ref{op1} asks.

\subsection*{Acknowledgement}
Computations were performed with Peanut, an open-source decision procedure for automatic
sequences written by the author; the argument was checked by independent brute-force
implementations.

\begin{thebibliography}{9}
\setlength{\itemsep}{0pt}\setlength{\parsep}{0pt}\setlength{\parskip}{0pt}

\bibitem{AS}
J.-P. Allouche and J. Shallit,
\emph{Automatic Sequences: Theory, Applications, Generalizations},
Cambridge University Press, 2003.

\bibitem{BSCRF}
F. Blanchet-Sadri, J. D. Currie, N. Rampersad, and N. Fox,
\emph{Abelian complexity of fixed point of morphism $0 \mapsto 012$, $1 \mapsto 02$,
$2 \mapsto 1$},
Electron. J. Combin. \textbf{14} (2014), \#A11.

\bibitem{Brlek}
S. Brlek,
\emph{Enumeration of factors in the Thue--Morse word},
Discrete Appl. Math. \textbf{24} (1989), 83--96.

\bibitem{Gelfond}
A. O. Gelfond,
\emph{Sur les nombres qui ont des propri\'et\'es additives et multiplicatives donn\'ees},
Acta Arith. \textbf{13} (1967/1968), 259--265.

\bibitem{Moradi-thesis}
D. Moradi,
\emph{State Complexity of Linear Relations and Linear Subsequences of Automatic Sequences},
MMath thesis, University of Waterloo, 2026. UWSpace,
\url{https://hdl.handle.net/10012/23044}.

\bibitem{MRS}
D. Moradi, N. Rampersad, and J. Shallit,
\emph{Complexity of linear subsequences of $k$-automatic sequences},
arXiv:2512.10017v5 [cs.FL], 2 April 2026;
also in Proc.\ CIAA 2026, LNCS \textbf{16695}, 185--199.

\bibitem{Starosta}
\v{S}. Starosta,
\emph{Generalized Thue--Morse words and palindromic richness},
Kybernetika \textbf{48} (2012), 361--370.

\bibitem{TrompShallit}
J. Tromp and J. Shallit,
\emph{Subword complexity of a generalized Thue--Morse word},
Inform. Process. Lett. \textbf{54} (1995), 313--316.

\end{thebibliography}

\end{document}
